\documentclass[10pt,a4paper]{amsart}
\usepackage[margin=19mm]{geometry}
\usepackage{amsmath,amssymb,mathtools}
\usepackage[scr=rsfs,cal=euler]{mathalfa}
\usepackage{xfrac}
\usepackage[unicode,hidelinks,bookmarks=false]{hyperref}
\newcommand{\ie}{i.e.\ }
\newcommand{\cf}{cf.\ }
\newcommand{\resp}{respectively\ }
\newcommand{\Subsection}{\subsection}
\providecommand{\nobreakdash}{\nobreak\hbox{-}}
\setlength{\emergencystretch}{2em}
\multlinegap0pt
\usepackage{mathtools}
\usepackage{amssymb}
\usepackage{bbm}
\usepackage{tikz}
\usepackage[shortlabels]{enumitem}
\usepackage{xcolor}
\usepackage{dsfont}
\usepackage[normalem]{ulem}
\newtheorem{lemma}{Lemma}
\newcommand{\of}[1]{\left( #1 \right)}
\newcommand{\set}[1]{\left\{ #1 \right\}}
\title{Chromatic polynomial evaluation spectra}
\author{Rafael Miyazaki, Cosmin Pohoata, Michael Zheng}
\date{Source: arXiv:2512.19600v1 (CC BY 4.0)}
\begin{document}
\maketitle

\noindent\emph{Source and licence.} Chromatic polynomial evaluation spectra. Version arXiv:2512.19600v1 (CC BY 4.0), distributed by arXiv under the Creative Commons Attribution 4.0 International licence, http://creativecommons.org/licenses/by/4.0/, as recorded in the arXiv licence metadata for this exact version and shown on the abstract page https://arxiv.org/abs/2512.19600v1. Full licence notice: \texttt{licenses/arxiv-2512.19600-CC-BY-4.0.txt}.\par\smallskip

\noindent\textit{Editorial scope.} Counted unit: the article's lemma lem:ping\_pong\_02 (source0.tex lines 890--894). The article's complete original proof of it is delivered with it (source0.tex lines 896--973, one \texttt{proof} environment of 78 lines). No mathematics is written, repaired or re-derived: the statement and the proof are the article's own lines, in the article's own notation, and no retained line is substituted or re-flowed.  Retained uncounted as the source-local context the proof needs: lines 728--748; lines 863--869.  Nothing was omitted from any retained span, so the package carries no omission record.  No retained line contains a citation, so the wrapper carries no bibliography and no bibliography entry.  No retained line contains a figure environment, an image-inclusion command or drawing code, so no image is delivered and the package declares no asset dependency.  Editorial formatting only, in addition: an AMS \texttt{amsart} preamble that restates the article's own declarations and macros used by the retained lines, namely the environments \texttt{lemma} on separate counters, together with the standard packages those lines need.  The retained environments print in this document as Lemma 1 in order of appearance, whereas the article prints the same items with the numbering of its own full document.  The complete native source of the article as distributed in the arXiv e-print for this exact version is bundled unchanged as \texttt{sources/191493/source0.tex}.
\begin{lemma}[Matrix action]\label{lem:matrices-pos}
Let $e$ be an edge of a planar graph $G$, and write $w(G,e)=\binom{t}{u}$.
Then
\[
(w\circ S)(G,e)
=
\begin{pmatrix}
-1 & 1\\
0 & q-1
\end{pmatrix}
\binom{t}{u},
\qquad
(w\circ B)(G,e)
=
\begin{pmatrix}
q-1 & 0\\
1 & q-2
\end{pmatrix}
\binom{t}{u}.
\]
\end{lemma}

\begin{lemma}
    \label{lem:S2_powers}
    For $m \geq 1$ and $q \neq 1$, we have 
    \begin{equation}\label{eq:ratio_S_2m}
        (r \circ w \circ S^{2m})(G,e) = \frac{1}{q} + \frac{(r \circ w)(G,e) - \frac{1}{q}}{(q-1)^{2m}}.
    \end{equation}
\end{lemma}

\begin{lemma}\label{lem:ping_pong_02} Let $q \in (0,2)\setminus \{1\}$. Then, there exists a planar graph $G_0$, an edge $e_0$ of $G_0$ and block operations $K=K(q)$ and $L=L(q)$ that are each a sequence of operations in $\{B,S\}$ such that the following is true.  If, for every word $a = (a_1,\dots,a_t)$ in the semigroup $\{K,L\}^\ast $, we let 
     \[\widetilde{\mathbf{w}}(a) \coloneqq (w\circ a_t\circ a_{t-1}\circ\dots\circ a_1)(G_0,e_0),\]
     the attainable vector witnessed by the pair generated through the sequence of operations encoded by $a$ starting with the pair $(G_0,e_0)$, then for all words $w\ne w'$ in the semigroup $\{K,L\}^\ast$ with the same number of letters,  $\widetilde{\mathbf{w}}(w)\neq \widetilde{\mathbf{w}}(w')$.
    
\end{lemma}

\begin{proof}


We start by remarking that in order to prove that a tuple $(G_0,e_0,K,L)$ satisfies the desired property, it is sufficient to find disjoint sets $I$ and $J$ such that \begin{enumerate}
    \item \label{item:1}$r(w(G_0,e_0)) \in I \cup J$;    \item\label{item:2}$r((w\circ K)(G,e))\in I$, for every edge $e$ of a planar graph $G$ such that $r(w(G,e)) \in I \cup J$;
    \item \label{item:3}$r((w\circ L)(G,e))\in J$, for every edge $e$ of a planar graph $G$ such that $r(w(G,e)) \in I \cup J$.
\end{enumerate}
Indeed, let $w, w' \in \set{K,L}^t$ for some $t \geq 0$ such that $\widetilde{\mathbf{w}}(w)= \widetilde{\mathbf{w}}(w')$. In particular, $(r \circ \widetilde{\mathbf{w}}) (w) = (r \circ \widetilde{\mathbf{w}})(w')$. If $ t= 0$, then we directly get $w = w'$. Otherwise, from \eqref{item:1}--\eqref{item:3} it follows that $(r\circ \widetilde{\mathbf{w}})(w'') \in I \cup J$ for any prefix $w''$ of $w$ or $w'$. Therefore, as $I$ and $J$ are disjoint, $w$ and $w'$ must agree in their last letter by \eqref{item:2} and \eqref{item:3}. As $q \not \in \set{0,1,2}$, $S$ and $B$ correspond to invertible matrices, hence the same is true for $K$ and $L$. Hence, we may invert to cancel the action of the last letter and iterate to conclude that $w = w'$.

We divide the proof into three cases, depending on $q$.

\emph{Case 1:} $q\in (0,3/2)$, $q\ne 1$. Let \[X =(G_0,e_0,K,L,I,J)=
\Big(K_4,e,S^{2m},B^2, \big(- \infty, (q-1)^2 / (2q-3)\big], \big((q-1)^2 / (2q-3),0\big)\Big),\]
where $e_0$ is any edge of $G_0=K_4$ and $m=m(q)$ is an integer to be chosen later.
We now prove that $X$ satisfies properties \eqref{item:1}--\eqref{item:3}.
\begin{enumerate}
    \item $r(w(G_0,e_0)) = 1/(q-2)\in (-\infty,0)=I \cup J$.
    \item By Lemma \ref{lem:S2_powers}, if $r = r(w(G,e)) \in (-\infty,0) = I \cup J$, then $w_k = (w \circ S^{2m})(G,e)$ satisfies
    \begin{equation*}
        r(w_k) = \frac{r}{(q-1)^{2m}} + \frac{1}{q} \of{1 - \frac{1}{(q-1)^{2m}}} < \frac{1}{q} \of{1 - \frac{1}{(q-1)^{2m}}} \overset{m \to \infty}{\longrightarrow} - \infty.
    \end{equation*}
       Hence, for sufficiently large $m$, we get $r(w_k) \in I$.  
    \item From Lemma~\ref{lem:matrices-pos}, we obtain
    \begin{align}\label{eq:B2_matrix}
        (w \circ B^2)(G,e) &= \begin{pmatrix}
            (q-1)^2 & 0 \\
            2q-3 & (q-2)^2 
        \end{pmatrix}
        w(G,e).
    \end{align}
    If $r = r(w(G,e)) \in (-\infty,0) = I \cup J$, then $w_\ell = (w \circ B^2)(G,e)$ satisfies 
    \begin{equation*}
        \frac{(q-1)^2}{2q-3}< r(w_\ell) = \frac{(q-1)^2}{(2q-3) + (q-2)^2 r^{-1}} < 0,
    \end{equation*}
    since $2q - 3 < 0$, $r<0$, and $q \neq 1$. Thus $r(w_\ell) \in J$.
\end{enumerate}
    

\emph{Case 2:} $q \in (3/2, 2)$. Let \[X =(G_0,e_0,K,L,I,J)=
\Big(K_3,e,S^{2m},B^2,  \big[(q-1)^2 / (2q-3),\infty\big), (1,(q-1)^2 / (2q-3))\Big),\]
where $e_0$ is any edge of $G_0=K_3$ and $m=m(q)$ is an integer to be chosen later.
We now prove that $X$ satisfies properties \eqref{item:1}--\eqref{item:3}.

\begin{enumerate}
    \item $r(w(G_0,e_0)) = 1/(q-1)\in (1,\infty)=I \cup J$.
    \item By Lemma~\ref{lem:S2_powers}, if $r = r(w(G,e)) \in (1,\infty) = I \cup J$, then $w_k = (w \circ S^{2m})(G,e)$ satisfies 
    \begin{equation*}
        r(w_k)= \frac{1}{q} + \frac{r - \frac{1}{q}}{(q-1)^{2m}} \geq \frac{1}{q} + \frac{1 - \frac{1}{q}}{(q-1)^{2m}} \overset{m \to \infty}{\longrightarrow} \infty
    \end{equation*}
    since $1> 1/q$. Hence, for sufficiently large $m$, we get $r(w_k) \in I$.
    \item By \eqref{eq:B2_matrix}, if $r = r(w(G,e)) \in (-\infty,0) = I \cup J$, then $w_\ell = (w \circ B^2)(G,e)$ satisfies
    \begin{equation*}
        1 = \frac{(q-1)^2}{(2q-3) + (q-2)^2}\leq r(w_\ell) = \frac{(q-1)^2}{(2q-3) + (q-2)^2 r^{-1}} < \frac{(q-1)^2}{2q-3}
    \end{equation*}
    since $q>3/2$ and $r>1$. Thus 
    $r(w_\ell)\in J$.
\end{enumerate}
\emph{Case 3:} $q = 3/2$. Let \[X =(G_0,e_0,K,L,I,J)=
\big(K_4,e,S^2,B\circ S^2\circ B, (-\infty,0], (0,1/2]\big),\]
where $e_0$ is any edge of $G_0=K_4$.
We now prove that $X$ satisfies properties \eqref{item:1}--\eqref{item:3}.
\begin{enumerate}
    \item $r(w(G_0,e_0)) = 1/(q-2) = -2\in (-\infty,1/2]=I \cup J$.
    \item By Lemma~\ref{lem:S2_powers}, if $r = r(w(G,e)) \in (-\infty,1/2] = I \cup J$, then \[w_k = (w \circ S^{2})(G,e) =4r-2 \in (-\infty,0] = I.\]
    
    \item From Lemma~\ref{lem:matrices-pos}, we obtain
     \begin{equation*}
        (w \circ L) (G,e) = \begin{pmatrix}
            0 & 1/8 \\
            -1/8 & 5/16
        \end{pmatrix} w(G,e). \end{equation*}
    Thus, if $r = r(w(G,e)) \in (-\infty,1/2]= I \cup J$, then $w_\ell = (w \circ L)(G,e)$ satisfies     \begin{equation*}
        0 < r(w_u) = \frac{2}{5 - 2 r } \leq \frac{1}{2},
    \end{equation*}
    since $r\mapsto 2/(5-2r)$ is strictly increasing in $I\cup J$. Thus 
    $r(w_\ell)\in J$. \qedhere
\end{enumerate}
\end{proof}

\end{document}
